% XeLaTeX can use any Mac OS X font. See the setromanfont command below.
% Input to XeLaTeX is full Unicode, so Unicode characters can be typed directly into the source.

% The next lines tell TeXShop to typeset with xelatex, and to open and save the source with Unicode encoding.

%!TEX TS-program = xelatex
%!TEX encoding = UTF-8 Unicode

\documentclass[11pt]{article}


%%%%%%%%%%%%%%%
\def\Type{ApplicationDJ}
\def\TypeNum{During-class for Application Day: Gradient Descent Method\\Key}
\def\TypeTitle{}
\def\Semester{Fall 2025}
\def\Course{Math 1190/1210}
\def\Targets{{\bf\color{Goldenrod}AD1}}
\def\pTableCellNum{9}
\def\learningRateCellNum{8}
\def\mbTableCellNum{11}
\def\exNum{$3$} %This exists because \label{} is not picking up $3$ from Example 3. Change if example ordering changes


\usepackage{preambleDJ}




\begin{document}
\pagestyle{otherpages}
\thispagestyle{firstpage}
\vspace*{.65in}

%If Ada Lovelace\footnote{\url{https://en.wikipedia.org/wiki/Ada_Lovelace}} is the ``mother of programming" then 
{\bf Introduction}\\
Fei-Fei Li is  known as the ``godmother of A.I" for her  groundbreaking contributions in computer vision and deep learning and as the creator of ImageNet\footnote{
See ``Godmother of A.I." Fei-Fei Li on technology development: ``The power lies within people" at \url{https://www.cbsnews.com/news/godmother-of-a-i-on-technology-development-fei-fei-li/}\\
 or Britannica at \url{https://www.britannica.com/biography/Fei-Fei-Li}, The Guardian \href{https://www.theguardian.com/technology/2023/nov/05/ai-pioneer-fei-fei-li-im-more-concerned-about- the-risks-that-are-here-and-now}{\underline{here}} or \url{https://en.wikipedia.org/wiki/Fei-Fei_Li} }. 
~\\\\
\mpageT{.35}{.65}{
%\graphicsR{adaLovelace2}{2}
\graphicsL{feiFeiLiWiredCropped}{2}
}
{
In our first application day involving the Keeling Curve, we modeled a subset of the CO$_2$ data (see chart below) using several lines of best fit. One of those was a linear function of the form $$c(t)=mt+b.$$  The process to find parameters $m$ and $b$ was done by Desmos behind the scenes.  Today we're going to use explore machine learning as one approach to estimate these parameters.    Note that $t=0$ corresponds to the year $1965$, $t=5$ is the year $1970$, etc.  Additionally the units for CO$_2$ are parts per {\em billion} (ppb).\footnotemark
\graphicsC{keelingCurve_1965-2025.pdf}{4}
}~\\
\footnotetext{This differs from Application Day: Keeling Curve where we used units of ``parts per million" }
Here's where we're headed:
\itemI{
\item We will interpret gradient vectors as a compass pointing in the direction of greatest ascent.
\item We will explore Gradient Ascent/Descent methods to find critical points related to local maximums/minimums.
\item We will use Gradient Descent method to minimize error/loss and estimate $m$ and $b$ from data.
}
\newpage
{\bf Revisiting Gradient Vectors}\\
We pick up where we left off in the pre-class activity with the surface $$\ds z=f(x,y)=\frac{1}{2}\left( x^3 +x^2y+xy^2-9x\right)$$ 
\enum{
\item Given the surface $\ds z=f(x,y)=\frac{1}{2}\left( x^3 +x^2y+xy^2-9x\right)$,  find $\ds |\nabla f(-1,3/2)|$.
\enum{
\item  Find $f_x(x,y)$. You may use your work from the pre-class activity. $f_x(x,y)=$\unline{1.5}{\blue{$\ds \frac{1}{2}\left( 3x^2+2xy +y^2-9 \right)$}}
\sol{
\al{
f_x(x,y)&=\frac{1}{2}\left( \frac{d}{dx}\left[x^3\right]+\frac{d}{dx}\left[ x^2y \right]+\frac{d}{dx}\left[xy^2\right]-\frac{d}{dx}\left[9x\right] \right)\\
&=\frac{1}{2}\left( \frac{d}{dx}\left[x^3\right]+y\frac{d}{dx}\left[ x^2 \right]+y^2\frac{d}{dx}\left[x\right]-\frac{d}{dx}\left[9x\right] \right)\qquad\text{\em \footnotesize We have moved constants   $y$ and $y^2$}\\
&=\frac{1}{2}\left( 3x^2+y\cdot 2x +y^2\cdot 1-9 \right)\\
&=\frac{1}{2}\left( 3x^2+2xy +y^2-9 \right)\\
}
}
\vs{-.4}
\item Find $f_y(x,y)$. You may use your work from the pre-class activity. $f_y(x,y)=$\unline{1.5}{\blue{$\ds \frac{1}{2}\left( x^2+2xy\right)$}}
\sol{
\al{
f_y(x,y)&=\frac{1}{2}\left( \frac{d}{dy}\left[x^3\right]+\frac{d}{dy}\left[ x^2y \right]+\frac{d}{dy}\left[xy^2\right]-\frac{d}{dy}\left[9x\right] \right)\\
f_y(x,y)&=\frac{1}{2}\left( \frac{d}{dy}\left[x^3\right]+x^2\frac{d}{dy}\left[ y \right]+x\frac{d}{dy}\left[y^2\right]-\frac{d}{dy}\left[9x\right] \right)\qquad\text{\em \footnotesize We have moved constants   $x$ and $x^2$}\\
f_y(x,y)&=\frac{1}{2}\left( 0+x^2 \cdot 1+x \cdot 2y-0 \right)\qquad\text{\em \footnotesize  $x^3$ is constant so $\ds \frac{d}{dy}[x^3]=0$. $9x$ is constant so $\ds \frac{d}{dy}[9x]=0$  }\\
f_y(x,y)&=\frac{1}{2}\left( x^2+2xy\right)
}
}
}
}
We evaluated $f_x(x,y)$ and $f_y(x,y)$ at $(x,y)=(-1,3/2)$ to get 
\itemI{
\item $\ds f_x(-1,3/2)=-\frac{27}{8}=-3.375$
\item $\ds f_y(-1,3/2)=-1$ 
\item Therefore our gradient vector was $\ds \nabla f(-1,3/2)=<f_x(-1,3/2),f_y(-1,3/2)>=<-3.375, -1>$
\item We also found the length of the gradient vector to be $$\ds |\nabla f(-1,3/2)|=|<-3.375, -1>|=\sqrt{(-3.375)^2+(-1)^2}\approx 3.52.$$
}


\newpage
{\bf Interpreting Gradient Vectors}\\
The gradient $\nabla f =<f_x,f_y>$ it has special properties related to the function $z=f(x,y)$.  One way to think of the gradient is that it acts as a ``compass". This compass-like behavior is going to help us locate critical points of local maximums and local minimums.  \\
\\
 {\bf \em Special Property of the Gradient Vector}
\itemI{
\item {\em At $(a,b)$, the gradient vector, $\nabla f(a,b)$ \blue{points in the direction our function $f$ is increasing the most at $(a,b,f(a,b))$}.}
\item {\em At $(a,b)$, the negative gradient vector, $-\nabla f(a,b)$ \blue{points in the direction our function $f$ is decreasing the most at $(a,b,f(a,b))$}.}
}


Let's visualize this vector and surface to better understand these properties.

\mpageT{.6}{.4}{
Go to \url{https://www.desmos.com/3d/9exv5rzkku}.  \\
Click on the wrench icon and check ``Reverse contrast", ``Translucent surfaces" and move the bounding box perspective slider to the far right for better viewing.
}{
\graphicsR{contrastTranslucentPerspective.pdf}{2.5}
}\\\\

{\bf \em Gradient Vector as ``compass"} \\
In Desmos, we see a default view is of our ``compass" in the $x/y$ plane.
\itemI{
\item The center of the compass is $P_0=(-1,3/2)$ (i.e. the green point). 
\item In Desmos, the $x$-coordinate for $P_0$ can also be seen in the table in cell $\#\pTableCellNum$ labeled $P.x$ The $y$-value of $P_0$ is labeled $P.y$
\item The gradient vector is the blue vector: $\nabla f(-1,3/2)=<-3.375,-1>$. In Desmos this is labeled as $V_{gradient}$
\item The negative gradient, $-\nabla f=<3.375,1>$, is the orange vector pointing in the opposite direction. 
\item In Desmos, $|\nabla f(-1,3/2)|$, the length of the gradient, is  labeled as $L_{engthGradient}$.
}
~\\
Toggle on/off the surface $\ds z=f(x,y)=\frac{1}{2}\left( x^3 +x^2y+xy^2-9x\right)$  by clicking on the gray rectangle.  You will see:\\
\mpageT{.65}{.5}{
\itemI{
\item A local maximum (i.e. a ``peak"). The red point in the $x/y$ plane is the critical point that lies directly below the local maximum.
\item  a local minimum (i.e. a ``valley"). The red point in the $x/y$ plane is the critical point that  lies directly above the local minimum.
\item The green  point $S_{current}=(-1,3/2,3.625)$ is on the surface directly above the base point of our gradient: $P_0=(-1,3/2)$.
}
}{
\graphicsC{localMaxMinLocationsLabeled}{2.5}
}
\itemI{
\item A copy of our gradient  (and negative gradient) vectors with base point  $S_{current}=(-1,3/2,3.625)$
}
{\bf \em ``Compass"}  If we were standing on the surface at the green point $S_{current}=(-1,3/2,3.625)$, the gradient $\nabla f=<-3.375,-1>$ (the blue vector) {\em points in the direction of greatest ascent}. In other words, {\em if we were to take a small step in that direction, we would initially encounter the steepest part of the hill.}
 
\enumR{
\item Click on the blue triangle to take one step. Our original base point for the gradient was $P_0=(-1,3/2)$. You will see a new green point in the $x/y$ plane that we will call $P_1$. {\em This is the new base point for the gradient vector.} $S_{current}$ has moved {\em farther} away from the local max and is no longer visible on our graph.  Note that a purple path connects the original point on the surface to our new point on the surface.
\enum{
\item Record the new point on the surface. Round to $3$ decimal places.\\ $S_{current}=(\unline{.75}{\blue{$-4.375$}} ,\unline{.75}{\blue{$0.5$}},\unline{.75}{\blue{$-17.944$}})$
\item Record the new gradient base point, $P_1$, in the $x/y$ plane. Round to $3$ decimal places.  In Desmos, you can find this value in the table in cell $\#\pTableCellNum$ for $n=1$. \\ $P_1=(\unline{.5}{\blue{$-4.375$}} ,\unline{.5}{\blue{$0.5$}})$
\item Record the length of the new gradient vector at $P_1$. Round to three decimal places (see $L_{engthGradient}$ in Desmos)\\
$|\nabla f(P_1)|$=\unline{.5}{\blue{$23.347$}}\\
{\em Note that this is much larger than $|\nabla f(P_0)|\approx 3.52$}
}
}
{\bf \em Step Size/Learning Rate}\\
Unfortunately, our first step took us {\em away from} the local maximum on the surface, {\em away from} the associated critical point in the $x/y$ with a {\em much larger} $L_{engthGradient}$.  We address this by shrinking the {\em step size} (also called the {\em learning rate}). 
\itemI{
\item Click on the blue circle to reset to $P_0=(-1,3/2)$ and $n=0$ steps.
\item Click on the yellow rectangle to toggle on/off a yellow vector. In Desmos, the default length of the yellow vector is $L=1$ (i.e. $100\%$) of the gradient vector $\nabla f(P_0)$. Here's the key idea:\\\\
{\em We will change our learning rate $L$ to a smaller size in order to take smaller steps in same direction as the gradient vector.}
}
\enumR{
\item In Desmos, change $L$ (i.e. the {\em learning rate} or {\em step-size}) from $L=1$ to $L=0.5$ (see cell $\#\learningRateCellNum$ in Desmos). This will ensure that each step will now be $50\%$ of the length of the gradient.  Take 5 steps.
\enum{
\item Record the last gradient base point, $P_5$, in the $x/y$ plane. Round to $3$ decimal places.  In Desmos, you can find this value in the table in cell $\#\pTableCellNum$ for $n=5$. \\ $P_5=(\unline{.5}{\blue{$-2.734$}} ,\unline{.5}{\blue{$1.343$}})$
\item Record the length of the gradient vector at $P_5$. Round to three decimal places (see $L_{engthGradient}$ in Desmos)\\
$|\nabla f(P_5)|$=\unline{.5}{\blue{$4.892$}}
}
{\em Note that while $|\nabla f(P_5)|$ is  larger than $|\nabla f(P_0)|\approx 3.52$, it is much better than what we saw when $L=1$. Additionally, all of our points on the surface are in the general vicinity of the local maximum. }
}
{\bf  Gradient Ascent Method}\\
The process of trying to find the critical point for a {\em local maximum} is known as the {\em Gradient Ascent Method}. We know critical points occur when both $f_x=0$ and $f_y=0$ and thus $|\nabla f|=0$.   Thus our goal with Gradient Ascent method is to:
\quot{Estimate the critical point of a local maximum by ...} \vs{-.2}
\quot{\hs{.2}... Creating a sequence of points in the $x/y$ plane via stepping in direction of the\\ \hs{.35} Gradient Vector and ...}\vs{-.3}
\quot{\hs{.2}... Stopping our sequence of points when $L_{engthGradient}=|\nabla f|\approx 0$.}

\enumR{
\item \label{q4} In Desmos, experiment with different values of $L$ so that $|\nabla f(P_n)|<0.05$ in $n=10$ steps or less. {\em Each time you choose a new $L$, re-set the algorithm  by clicking the blue circle.}
\enum{
\item Record $L=$\unline{.5}{\blue{$0.3$}} \blue{There are many values that could work. $L=0.3$ is one such example.}
\item Record number of steps needed to achieve $|\nabla f(P_n)|<0.05$ (i.e. $L_{engthGradient}<0.05$)\\
$n=$\unline{.5}{\blue{$5$}} \blue{This number will depend on choice of $L$}
\item Record approximate critical point at final step $P_n$. Round to $3$ decimal places.  (see table for $P.x$ and $P.y$ in cell $\#\pTableCellNum$).\\
$P_n=(\unline{.5}{\blue{$-1.999$} }, \unline{.5}{\blue{$1.046$} })$ \blue{This point will depend on choice of $L$. Note that actual critical point is $(-2,1)$.}
\item Record $|\nabla f(P_n)|$ (i.e. $L_{engthGradient}$). Round to $3$ decimal places.\\
$|\nabla f(P_n)|=\unlineb{.5}{0.031}$ \blue{This number will depend on choice of $L$}
\item Record the local maximum value. Round to 3 decimal places (see $z$ coordinate of $S_{current}$).
$z=f(P_n)=$\unlineb{.5}{6.000}
}
}
The sequence of (green) points generated by Gradient Ascent method to locate the local maximum critical point can be seen in the $x/y$ plane.\\\\
{\bf \em Gradient Ascent Summary}\\
What we've seen so far is an algorithm with the following pattern
\quot{Start at a base point $P_0$, Find the gradient vector and step $L\%$ along the gradient vector to find a new $P_1$ that is closer to the critical point. Repeat until ``close enough".}
Mathematical notation that says the same thing is
\al{
P_0 +L\cdot \nabla f(P_0)=P_1
\intertext{More generally, if we have already taken $n$ steps to arrive at $P_n$, we can find $P_{n+1}$ by}
P_n +L\cdot \nabla f(P_n)=P_{n+1}
}

\example When $P_0=(-1,3/2)$,  we've seen $\nabla f(P_0)=<-3.375,-1>$.  Then if $L=0.5$, 
\al{
P_1&=P_0+L \cdot \nabla f(P_0)\\
&=(-1,1.5) + 0.5\cdot <-3.375,\,-1>\\
&=(-1,1.5) + <(0.5)\cdot (-3.375),\,(0.5)\cdot(-1)>\\
&=(-1,1.5) +<-1.6875 , \, -0.5>\\
&= (-1+-1.6875,\,1.5-.5)\\
P_1&=(-2.6875, 1)
}
Thus our first estimate of the (red) critical point in the $x/y$ plane is  $P_1=(-2.6875, 1)$ with  local maximum value of $f(P_1)=f(-2.6875, 1)\approx 4.656$. In question $\#\ref{q4}$, Desmos refined this estimate by taking several more steps until $|\nabla f(P_n))|=L_{engthGradient}<0.05$.
\\\\
{\bf  Gradient Descent Method}\\
The process of trying to find the critical point for a {\em local minimum} is known as the Gradient {\em  Descent} Method.  The primary difference between Gradient Descent and Gradient Ascent is that we step in the direction the function $f$ is {\em decreasing} the fastest. This will be in the  direction of the {\em negative} gradient vector. That means our goal is:
\quot{Estimate the critical point of a local minimum by ...} 
\quot{\hs{.2}... Creating a sequence of points in the $x/y$ plane via stepping in {\bf negative} direction of the Gradient Vector and ...}
\quot{\hs{.2}... Stopping our sequence of points when $L_{engthGradient}=|\nabla f|\approx 0$.}
Mathematically, if our starting base point is $P_0$, our algorithm becomes
\al{
P_0 -L\cdot \nabla f(P_0)=P_1
\intertext{More generally, if we have already taken $n$ steps to arrive at $P_n$, we can find $P_{n+1}$ by}
P_n -L\cdot \nabla f(P_n)=P_{n+1}
}
\enumR{
\item In Desmos, start over by clicking the blue circle to reset to $P_0=(-1,3/2)$. Locate the {\em local minimum}  by experimenting with different values of $L$ so that $|\nabla f(P_n)|<0.05$ in less than $n=15$ steps. {\em Each time you choose a new $L$, start over by clicking the blue circle.} {\bf In order to step in the direction of the negative gradient  (orange) vector, click on the orange triangle}.
\enum{
\item Record $L=$\unlineb{.5}{0.32} \blue{There are many values that could work. $L=0.32$ is one such example.}
\item Record number of steps needed to achieve $|\nabla f(P_n)|<0.05$ (i.e. $L_{engthGradient}<0.05$)\\
$n=$\unlineb{.5}{12} \blue{This number will depend on choice of $L$}
\item Record approximate critical point at final step $P_n$. Round to $3$ decimal places.  (see table for $P.x$ and $P.y$ in cell $\#\pTableCellNum$).\\
$P_n=(\unlineb{.5}{ 2.006}, \unlineb{.5}{ -.990})$ \blue{This point will depend on choice of $L$. The actual critical point is $(2,-1)$.}
\item Record $|\nabla f(P_n)|$ (i.e. $L_{engthGradient}$). Round to $3$ decimal places.\\
$|\nabla f(P_n)|=\unlineb{.5}{0.047}$ \blue{This number will depend on choice of $L$}
\item Record the local minimum value. Round to 3 decimal places. (See $z$ coordinate of $S_{current}$
$z=f(P_n)=$\unlineb{.5}{-6.000}
}
}
The sequence of (green) points generated by Gradient Descent method to locate the local minimum critical point can be seen in the $x/y$ plane.
\enumR{
\item {\bf (Extra Practice, if time)} Given $f(x,y)=x^2-x+y+y^2$, we will estimate the location of the local minimum using the Gradient Descent method using $P_0=(2,3)$ and a step size of $\ds L=\frac{1}{4}$.
\enum{
\item Find $f_x(x,y)$ and $f_y(x,y)$.
\sol{
\al{
f_x(x,y)&=\frac{d}{dx}[x^2]-\frac{d}{dx}[x]+\frac{d}{dx}[y]+\frac{d}{dx}[y^2]\\
&=2x-1+0 +0 \qquad \text{ (since $y$ and $y^2$ are constants relative to $x$)}\\
&=2x-1\\
f_y(x,y)&=\frac{d}{dy}[x^2]-\frac{d}{dy}[x]+\frac{d}{dy}[y]+\frac{d}{dy}[y^2]\\
&=0-0+1+2y  \qquad \text{ (since $x^2$ and $x$ are constants relative to $y$)}\\
&=1+2y
}
}
\item Find $\nabla f(P_0)=\nabla f(2,3)$
\sol{
\al{
f_x(2,3)&=2\cdot 2-1=3\\
f_y(2,3)&=1+2\cdot 3=7
\intertext{thus,}
\nabla f(2,3)&=<f_x(2,3), f_y(2,3)>\\
&=<3,7>
}
}
\item Find $P_1$.  Since we are searching for a local minimum, we use Gradient Descent:\\ $P_1=P_0-L\cdot \nabla f(P_0)$
\sol{
\al{
P_1&=P_0-L\cdot \nabla f(P_0)\\
&=(2,3)-\frac{1}{4}\cdot <3,7>\\
&=(2,3)-\left< \frac{3}{4},\frac{7}{4} \right>\\
&=\left(2-\frac{3}{4},3-\frac{7}{4} \right)\\
&=\left( \frac{5}{4}, \frac{5}{4}\right)
}
}
\item Find $P_2$.  Since we are searching for a local minimum, we use Gradient Descent:\\  $P_2=P_1+L\cdot \nabla f(P_1)$. Technology might be helpful.
\sol{
\al{
P_2&=P_1-L\cdot \nabla f(P_1)\\
&=\left( \frac{5}{4}, \frac{5}{4}\right)-\frac{1}{4}\nabla f(5/4,5/4)
\intertext{$\ds f_x(5/4,5/4)=\frac{3}{2}$ and $\ds f_y(5/4,5/4)=\frac{7}{2}$ thus,}
&=\left( \frac{5}{4}, \frac{5}{4}\right)-\frac{1}{4} \left<\frac{3}{2},\frac{7}{2}  \right>\\
&=\left( \frac{5}{4}, \frac{5}{4}\right)- \left<\frac{3}{8},\frac{7}{8}  \right>\\
&=\left( \frac{7}{8}, \frac{3}{8}\right)
}
}
}
More iterations returns $P_{10}\approx(.501,-.497)$ which is somewhat close\footnote{What ``close enough" means is an important topic that will depend on the context for using Gradient Ascent/Descent} to the actual minimum at $(1/2,-1/2)$.
}~\\
\newpage
{\bf Minimizing Total Error}\\
Recall that our goal today is to estimate parameters $m$ and $b$ so that we can model our CO$_2$ data with a linear function of the form $$c(t)=mt+b.$$  

\mpageT{.65}{.35}[\hs{.2}]{
Since we are estimating $m$ and $b$, we hope to minimize the error of our estimate. Error (sometimes called {\em loss}) is  how ``close" our proposed linear function $c(t)=mt+b$ gets to  data points.  We  measure the error/loss by finding the difference between the CO$_2$ value from data, $c_i$, and the predicted CO$_2$ value, $c(t_i)=mt_i+b$.  Lastly we square this difference to guarantee it will be positive. %
}{
\graphicsCFig{errorLoss1Point}{2.5}{Error (or loss) between  data point and function $c(t)$. \label{fig1}}
}\\
\example In Figure \ref{fig1}, we have a single data point $(t_1,c_1)=(5,0.325)$  We measure the error, $e_1$, between it and the model  $c(t)=mt+b$ by finding the difference.
\al{
e_1&=c_1-c(t_1) \qquad \text{\small [difference between amount of CO$_2$ from data, $c_1$, and amount predicted by the function, $c(t_1)$]}\\
e_1&=0.325-(5m+b) \qquad\text{[\small $c_1=0.325$ and $t_1=5$]}\\
e_1&=0.325-5m-b
}
We then square the error to make it positive:
$$e^2_1=(0.325-5m-b)^2$$
{\bf \em Total Error/Loss Function}
The {\em total} error/loss function is the average of the errors for  $n$ data points. $$\mathcal{L} (m,b)=\frac{1}{n}(e^2_1+e^2_2+e^2_3+e^2_4+e^2_5+\cdots +e^2_n)$$ Note that this is a function of $m$ and $b$. 
\begin{center}{\em $\mathcal{L} (m,b)$ is the function we will minimize  by using the Gradient Descent Method}\end{center}
\mpageT{.6}{.4}{
\example  \label{ex3}  Let's use $n=2$ points from our CO$_2$ data, $P_1=(5,.325)$, %$P_2=(20,.345)$ 
and $P_2=(40,.379)$,  to illustrate how this works (See Figure \ref{fig2}).
\al{
e^2_1&=(.325-c(5))^2=(.325-(5m+b))^2=(.325-5m-b)^2\\
%e^2_2&=(.345-c(20))^2=(.345-(20m+b))^2=(.345-20m-b)^2\\
e^2_2&=(.379-c(40))^2=(.379-(40m+b))^2=(.379-40m-b)^2
}
}{
\graphicsCFig{errorLoss2Points}{2.5}{Example error/loss. \label{fig2}}
}

\mpageT{.6}{.4}{
Thus,
\al{
\mathcal{L}(m,b)&=\frac{1}{2}(e^2_1+e^2_2)\\
\mathcal{L}(m,b)&=\frac{1}{2}((.325-5m-b)^2+(.379-40m-b)^2)
}
This is a function of two variables, $m$ and $b$, which can be minimized using the Gradient Descent method. See graph of $\mathcal{L}(m,b)$ in Figure \ref{fig3}.
}{
\graphicsCFig{errorLossFunction2Points}{2.5}{Graph of $\mathcal{L}(m.b)$. \label{fig3}}
}
 

\enumR{
\item Find the gradient vector for the error/loss function  $\mathcal{L}(m,b)$ in  Example \exNum.
$$\mathcal{L}(m,b)=\frac{1}{2}((.325-5m-b)^2+(.379-40m-b)^2)$$
Your answer will contain variables $m$ and $b$.
\enum{
\item Find $\mathcal{L}_m(m,b)$.  Your answer will contain variables $m$ and $b$.\\
$\mathcal{L}_m(m,b)=\unlineb{1.5}{}$ 
\sol{
\al{
\frac{d}{dm}\Bigg[\frac{1}{2} \Big((.325&-5m-b)^2+(.379-40m-b)^2\Big)\Bigg]\hs{.15}\text{\footnotesize (Use the Chain Rule on each term with $b$ constant)}\\
&=\frac{d}{dm}[\frac{1}{2}(.325-5m-b)^2]+\frac{d}{dm}[\frac{1}{2}(.379-40m-b)^2]\\
&=\frac{1}{2}\cdot 2\cdot(.325-5m-b)(-5)+\frac{1}{2}\cdot2\cdot(.379-40m-b)(-40) \hs{.15}\text{\footnotesize (It's ok to stop here}\\
&\, \hs{4.5} \text{\footnotesize and not simplify further) }\\
&=-5\cdot(.325-5m-b)-40\cdot(.379-40m-b)\\
&=-1.625+25m+5b-15.6+1600m+40b\\
\mathcal{L}_m(m,b)&=-16.785+1625m+45b
}
}\vs{-.2}
\item Find $\mathcal{L}_b(m,b)$.  Your answer will contain variables $m$ and $b$.\\
$\mathcal{L}_b(m,b)=\unlineb{1.5}{}$ 
\sol{
\al{
\frac{d}{db}\Bigg[\frac{1}{2}\Big(.325&-5m-b)^2+(.379-40m-b)^2\Big)\Bigg]\hs{.15}\text{\footnotesize (Use the Chain Rule on each term with $m$ constant)}\\
&=\frac{d}{db}[\frac{1}{2}(.325-5m-b)^2]+\frac{d}{db}[\frac{1}{2}(.379-40m-b)^2]\\
&=\frac{1}{2}\cdot 2\cdot(.325-5m-b)(-1)+\frac{1}{2}\cdot 2\cdot(.379-40m-b)(-1)\hs{.15}\text{\footnotesize (It's ok to stop here}\\
&\, \hs{4.5} \text{\footnotesize and not simplify further) }\\
&=-(.325-5m-b)-(.379-40m-b)\\
&=-.325+5m+b-.379+40m+b\\
\mathcal{L}_b(m,b)&=-0.704+45m+2b
}
}
}
}
 \vs{-.35}
{\bf Machine Learning and the  Gradient Descent Method}\\
We are finally ready to estimate $m$ and $b$ for use in our linear model $$c(t)=mt+b.$$
From here, Desmos will use the Gradient Descent method to minimize the total error. This will create a sequence of points of the form $P=(m.b)$ that will converge on the critical point associated with the minimum of our error/loss function $\mathcal{L} (m,b)$.
The data we're going to use can be found \href{https://www.desmos.com/calculator/npjnjtbrbz}.

{\underline{here}}.
We're going to locate the local minimum in three stages.
\itemI{
\item[Stage $1$:] The {\em training} stage utilizes the first $80\%$ of the data to produce a rough estimate of the critical point $(m,b)$
\item[Stage $2$:]  The {\em refinement and validation} stage on the next $10\%$ of the data to refines our estimate of the critical point $(m,b)$
\item[Stage $3$:] The {\em test} stage on the last $10\%$ of the data to determines if our linear mode $c(t)=mt+b$ is an effective model or not.
}
\mpageT{.6}{.4}{
Go to \url{https://www.desmos.com/3d/r4neljvinh}.  \\
Click on the wrench icon and check ``Reverse contrast", ``Translucent surfaces" and move the bounding box perspective slider to the far right for better viewing.
}{
\graphicsR{contrastTranslucentPerspective.pdf}{2.5}
}\\

In Desmos, you will see:
\itemI{
\item A graph of the error/loss function for the first $80\%$ of the data that looks like a very skinny paraboloid. In particular this uses the first $48$ data points and 
$$\mathcal{L} (m,b)=\frac{1}{48}\left(e^2_1+e^2_2+e^2_3+e^2_4+e^2_5+\cdots +e^2_{48}\right).$$
\item The orange rectangle advances the Gradient Descent several steps at one time.  It stops when one of the following stopping criteria is met.
\itemI{
\item The total error $\mathcal{L}(m,b)$  is below an error threshold of $0.001$. In Desmos, the total error on the training data is labeled $E_{rrorTrain}$ and the error threshold is labeled $E_{rrorThreshold}=0.001$.  
\item $|\nabla \mathcal{L}(m,b)|>200$ (i.e. an upper threshold). In Desmos this is labeled $L_{engthGradient}>200$. The upper threshold detects a situation when our estimates are moving {\em away} from the desired critical point and shuts down the algorithm.
}
\item The blue circle re-sets updates all new values and resets to $P_0=(m_{start},b_{start})$
\item A red rectangle  appears when algorithm is running (i.e. orange has been clicked) and will pause Gradient Descent.
\item $m_{start}=1$ is our initial slope. $b_{start}=1$ is our initial vertical intercept. 
\item $L=.001$ is the initial {\em learning rate} or {\em step size}
\item The default error threshold we will use to decide if we are ``close enough" to the critical point is $E_{rrorThreshold}=0.001$
\item Total Error on the training stage is $\mathcal{L}_{train}(m,b)$. In Desmos this is labeled $E_{rrorTrain}$
\item $|\nabla \mathcal{L}(m,b)|$ is labeled $L_{engthGradient}$ in Desmos.
\item Total Error on the validation stage is $\mathcal{L}_{validate}(m,b)$. In Desmos this is labeled $E_{rrorValidate}$
\item Total Error on the test stage is $\mathcal{L}_{test}(m,b)$. In Desmos this is labeled $E_{rrorTest}$
\item The sequence of $(m,b)$ points from the Gradient Descent method can be found in the table in cell $\#\mbTableCellNum$
}
{\bf \em Training Stage}
In the following exercise we use the training data (i.e. first $80\%$) to make our first estimate of $m$ and $b$.  This requires balancing accuracy of our estimation with computational cost. The more iterations are need, the more ``expensive" the computation. But fewer iterations may result in less accuracy.
\enumR{
\item Find a suitable value for $(m_{start},b_{start})$ that will converge within $250$ iterations or less and $E_{rrorTrain}<E_{rrorThreshold}=0.001$  \\\\
Note that the default value starting point  $P_0=(m_{start},b_{start})=(1,1)$ will not converge to the critical point because the stopping criteria $|\nabla \mathcal{L}| =L_{engthGradient}=1518.95>200$ is already met. 
\itemI{
\item You may find the Keeling Curve data useful to estimate $(m_{start},b_{start})$: \url{https://www.desmos.com/calculator/npjnjtbrbz}
\item After changing $(m_{start},b_{start})$, click the blue circle to update values in Desmos. Click Orange rectangle to start Gradient Descent algorithm.
\item When algorithm stops, check to make sure that it is due to $E_{rrorTrain}<E_{rrorThreshold}=0.001$ and {\em not} due to the fact that $L_{engthGradient}>200$. If it is the latter, you will need to try  different starting values for  $(m_{start},b_{start})$.
}
\enum{
\item Record $(m_{start},b_{start})=(\unlineb{.5}{0},\unlineb{.5}{0.3})$ \blue{There are many reasonable answers. Most will require $b_{start}$ to be somewhere around $0.32$.}
\item Record the number of iterations needed to converge. See table in Desmos, cell $\#\mbTableCellNum$\\
 $n=\unlineb{.5}{1}$  \blue{This will depend on choice of $(m_{start},b_{start})$.}
\item Record the final estimate of $(m,b)$ Round to $4$ decimal places. See table in Desmos, cell $\#\mbTableCellNum$\\
 $(m,b)=(\unlineb{.5}{0.0031}, \unlineb{.5}{.3001})$ \blue{This will depend on choice of $(m_{start},b_{start})$.}
}
}
~\\
{\bf \em Refining and validating the final estimate of $(m,b)$} We now a) refine our initial estimate and then b) validate it by ensuring the total error on the validation data (i.e.  next $10\%$ of original data set) is also within our error threshold.  
\enumR{
\item \label{refine} {\em Refining} Refine our final estimate of $(m,b)$ by reducing the threshold to\\ $E{rrorThreshold}=0.00001$ (four zeroes and a $1$ after the decimal). You will likely need to update the learning rate $L$ and/or the starting point $(m_{start},b_{start})$ to converge to the critical point in less than $250$ iterations and $E_{rrorTrain}<E_{rrorThreshold}=0.00001$.
\itemI{
\item After changing $L$ or $(m_{start},b_{start})$, click the blue circle to update values in Desmos and then the Orange rectangle to start Gradient Descent algorithm.
\item When algorithm stops, check to make sure that it is due to $E_{rrorTrain}<E_{rrorThreshold}=0.001$. If it is due to the fact that $L_{engthGradient}>200$, you will need to try  different starting values for  $(m_{start},b_{start})$.
}
\enum{
\item Record $L=\unlineb{.5}{0.001}$ \blue{There are many combinations of $L$ and $(m_{start},b_{start})$ that will converge.  } 
\item Record $(m_{start},b_{start})=(\unlineb{.5}{0},\unlineb{.5}{0.32})$ \blue{I left the default $L=0.001$ and changed $b_{start}$ from $0.3$ to $0.32$}
\item Record the number of iterations needed to converge. See table in Desmos, cell $\#\mbTableCellNum$\\
 $n=\unlineb{.5}{5}$  \blue{This will depend on choice of $(m_{start},b_{start})$ and $L$.}
\item Record the final estimate of $(m,b)$ Round to $4$ decimal places. See table in Desmos, cell $\#\mbTableCellNum$\\
 $(m,b)=(\unlineb{.5}{0.0015}, \unlineb{.5}{.3201})$ \blue{This will depend on choice of $(m_{start},b_{start})$ and $L$.}
}
\item Compare validation error close to the error threshold. Is  $E_{rrorValidate}$ within $ 0.00008$ of $E_{rrorThreshold}=0.00001$? That is, if $E_{rrorValidate}$ is in the interval $[0,0.00009]$, then we're going to say "close enough" and record values below. If not, you will need to repeat your experiment in question \ref{refine}. When $E_{rrorValidate}$ within $ 0.00008$ of $E_{rrorThreshold}$, record values below.\enum{
\item Record $L=\unlineb{.5}{0.001}$ \blue{There are many combinations of $L$ and $(m_{start},b_{start})$ that will converge.  } 
\item Record $(m_{start},b_{start})=(\unlineb{.5}{0},\unlineb{.5}{0.32})$ \blue{I left the default $L=0.001$ and changed $b_{start}$ from $0.3$ to $0.32$}
\item Record the number of iterations needed to converge. See table in Desmos, cell $\#\mbTableCellNum$\\
 $n=\unlineb{.5}{5}$  \blue{This will depend on choice of $(m_{start},b_{start})$ and$L$.}
\item Record the final estimate of $(m,b)$ Round to $4$ decimal places. See table in Desmos, cell $\#\mbTableCellNum$\\
 $(m,b)=(\unlineb{.5}{0.0015}, \unlineb{.5}{.3201})$ \blue{This will depend on choice of $(m_{start},b_{start})$ and $L$.}
 \item Record $E_{rrorValidate}$ up to 6 decimal places (several decimal places will be $0$)\\
 $E_{rrorValidate}=\unlineb{.75}{0.000078}$
}
}
{\bf \em Testing stage} Our last step is to examine the total error using the Test data (i.e. last $10\%$ of our original data set. 
\itemI{
\item Is $E_{rrorTest}$ within $ 0.00008$ of $E_{rrorThreshold}=0.00001$? That is, if $E_{rrorTest}$ is in the interval $[0,0.00009]$, then we're going to say ``close enough" and conclude we have an effective model.
\item  If   $E_{rrorTest}$ is {\em not} in the interval $[0,0.00009]$ then we will conclude that our linear model $c(t)=mt+b$ is probably not the best model to use for this data set and we should consider another model such as an exponential model.
}
\enumR{
\item Examine $E_{rrorTest}$ in Desmos.\enum{
\item Record $E_{rrorTest}$ up to $6$ decimal places.  $E_{rrorTest}=\unlineb{.75}{0.000238}$\\
\item Using our estimated values for $(m,t)$, Is $c(t)=mt+b$ a good model of our data set? Hint: examine whether or not $E_{rrorTest}$ in the interval $[0,0.00009]$. \\
 \blue{No, $E_{rrorTest}$ is too large. It is likely that an exponential model would do a better job. If we take our final values of $m$ and $b$ and use them to graph $c(t)=mt+b$ with our \href{https://www.desmos.com/calculator/npjnjtbrbz}{\underline{Keeling Curve data set}}, we see the last part of our data set pulls away from the line. This likely explains the greater error. }
}

}




\end{document}  
